% !TeX spellcheck = ru_RU
% !TEX root = vkr.tex

\section{Тестирование, верификация и методическое обеспечение}
\label{sec:testing}

В разделе представлена методика проверки разработанного программного комплекса, результаты функциональной верификации его компонентов (процессорное ядро RISC‑V RV32I, моделируемая платформа, операционная система xv6, а также пользовательская утилита \texttt{tinyc} на основе компилятора TinyCC, встроенная в образ ОС), а также описание учебно‑методических материалов, предназначенных для использования в курсах по архитектуре ЭВМ и системному программированию.

\subsection{Методика тестирования и конфигурация стенда}

Тестирование организовано на трёх уровнях, соответствующих иерархии компонентов: верификация системы команд RV32I, проверка выполнения скомпилированных C‑программ на платформе без операционной системы (режим \texttt{run\_c\_version}) и функциональная верификация операционной системы xv6 (режим \texttt{xv6\_version}). Для каждого уровня определены тестовые сценарии, критерии прохождения и используемые инструменты.

\textbf{Конфигурация стенда.} Симуляция Verilog‑моделей выполняется на рабочих станциях с процессорами Intel Core i3/i7 (Ubuntu 22.04). Воспроизводимость окружения обеспечена Docker‑образом (на основе \texttt{ubuntu:22.04}), включающим Icarus Verilog, Verilator, GTKWave, кросс‑тулчейн \texttt{riscv64-unknown-elf-gcc} (12.2.0), Python 3 с вспомогательными скриптами (\texttt{encode.py}, \texttt{elf2hex.py}). Все тесты запускаются через корневой Makefile, автоматизирующий сборку, симуляцию и сравнение выходных логов.

\textbf{Методика верификации процессорного ядра.} Разработан набор ассемблерных тестов, охватывающий все форматы инструкций RV32I. Каждый тест загружается в память инструкций через \texttt{\$readmemh}. Корректность выполнения оценивается двумя способами:
\begin{itemize}
\item сравнение итогового состояния регистрового файла и памяти данных с эталоном, полученным из независимого C++ ISA‑симулятора;
\item анализ потактового лога и VCD‑диаграмм (GTKWave) для визуального контроля временных соотношений.
\end{itemize}
Тест считается пройденным при полном совпадении всех регистров (кроме x0) и памяти данных, а также при отсутствии неопределённых состояний в сигналах.

\textbf{Методика тестирования платформы в режиме \texttt{run\_c\_version}.} Проверяется выполнение скомпилированного C‑кода без операционной системы. Тестовая программа вычисляет факториал и числа Фибоначчи (рекурсивная и итеративная версии). Она компилируется кросс‑компилятором \texttt{riscv64-unknown-elf-gcc}, линкуется с минимальной библиотекой поддержки (точка входа \texttt{\_start}, инициализация стека, \texttt{putchar}/\texttt{getchar}) и преобразуется в hex‑образ. Критерий прохождения – вывод вычисленных значений через эмулируемый UART в файл \texttt{uart\_output.txt}, совпадающий с эталонной строкой.

\textbf{Методика верификации платформы в режиме \texttt{xv6\_version}.} Включает проверку загрузки ядра, выполнения системных вызовов, многозадачности и работы файловой системы на RAM‑диске. Сценарии:
\begin{enumerate}
\item \textbf{Загрузка ядра} – после сброса процессор начинает выполнение с адреса 0x80000000, инициализируются устройства, создаётся первый процесс (\texttt{init}) и запускается оболочка (\texttt{sh}). Критерий – появление приглашения shell в логе UART.
\item \textbf{Выполнение встроенных команд оболочки} (\texttt{ls}, \texttt{cat}, \texttt{echo}, \texttt{mkdir}, \texttt{rm}). Ожидаемый результат – вывод соответствующих сообщений, отсутствие паники ядра.
\item \textbf{Многозадачность} – запуск двух фоновых процессов, проверка чередования вывода. Анализ VCD‑диаграмм подтверждает переключение контекста по прерываниям таймера.
\item \textbf{Системные вызовы} – выполнение пользовательской программы с \texttt{fork}, \texttt{exec}, \texttt{wait}, \texttt{read}, \texttt{write}. Критерий – корректное создание и завершение дочернего процесса без утечек памяти.
\end{enumerate}
Для ускорения отладки на ранних этапах используется C++ ISA‑симулятор, финальная верификация выполняется в RTL‑симуляции Verilator.

\subsection{Результаты функционального тестирования и верификации}

\textbf{Процессорное ядро RV32I.} Все инструкции базового целочисленного набора прошли проверку на 150 тестовых случаях, включая граничные значения, знаковое расширение, выравнивание адресов и генерацию исключений. Состояние регистров и памяти после выполнения каждого теста совпало с ISA‑симулятором. Временные диаграммы подтвердили, что комбинационные пути укладываются в один такт при частоте 100 МГц.

\textbf{Платформа \texttt{run\_c\_version}.} Тестовые программы на C (факториал 5, числа Фибоначчи до 10) скомпилированы штатным кросс‑компилятором \texttt{riscv64-unknown-elf-gcc} и успешно выполнены на виртуальном процессоре в режиме \texttt{run\_c\_version}. Выходные строки в \texttt{uart\_output.txt} совпали с ожидаемыми: \texttt{factorial(5)=120}, \texttt{fib(10)=55}. Это подтверждает корректность подсистемы памяти, инструкций загрузки/сохранения и интерфейса с UART.

\textbf{Платформа \texttt{xv6\_version}.} Получены следующие результаты:
\begin{itemize}
\item Загрузка ядра: в логе UART появилось приглашение оболочки, инициализация драйверов и базовых подсистем (включая настройку физической памяти в M‑mode) завершена без ошибок.
\item Команды оболочки: \texttt{ls}, \texttt{cat}, \texttt{echo}, \texttt{mkdir}, \texttt{rm} отработали штатно; содержимое файлов совпало с эталоном.
\item Многозадачность: при запуске двух фоновых процессов вывод чередовался, планировщик round‑robin переключал контекст каждые 10 мс. VCD‑диаграмма показала корректное сохранение регистров и переключение стека по прерыванию таймера.
\item Системные вызовы: программа с \texttt{fork} и \texttt{exec} создала дочерний процесс, выполнила \texttt{/bin/echo} и завершилась. Утилита \texttt{memtest} не выявила утечек памяти.
\end{itemize}
Портированный TinyCC внутри xv6 успешно скомпилировал и выполнил программу вычисления чисел Фибоначчи, что демонстрирует цикл разработки на целевой платформе.

Все функциональные тесты завершены без ошибок; ядро xv6 не паниковало, пользовательские программы не вызывали исключений. Работоспособность комплекса подтверждена.

\subsection{Учебно‑методические материалы}

На основе разработанного комплекса подготовлен набор материалов для курсов «Архитектура ЭВМ» и «Системное программирование»:

\begin{itemize}
\item Лекционная презентация «От транзистора до операционной системы», охватывающая набор данных RISC‑V, систему команд, механизмы исключений, управление памятью и процессами на примере xv6, а также роль компилятора в жизненном цикле программы.
\item Документация по архитектуре комплекса, включающая: компонентную структуру, карту адресного пространства (RAM, UART, CLINT, RAM‑диск), инструкции по сборке и запуску для двух режимов, руководство по написанию ассемблерных программ (с использованием \texttt{encode.py}), описание скриптов верификации.
\item Сценарии симуляции (Makefile‑цели и тестбенчи) для лекционных демонстраций: запуск ассемблерной программы с выводом регистров, выполнение C‑программы «Hello, world» на «голой» платформе, загрузка xv6 с командами \texttt{ls}, \texttt{cat}, \texttt{tinyc}, генерация VCD‑файла и анализ в GTKWave для визуализации переключения контекста.
\item Комплект учебных заданий, каждое из которых содержит цель, теоретическую справку, порядок выполнения и контрольные вопросы: знакомство с RISC‑V ассемблером и \texttt{encode.py}; модификация процессорного ядра (добавление новой инструкции); реализация драйвера UART на C для режима без ОС; добавление нового системного вызова в xv6; создание пользовательской программы, скомпилированной TinyCC внутри xv6.
\item Рекомендации по педагогической оценке (анкеты для студентов и преподавателей) для сбора обратной связи о степени понимания материала и предложений по улучшению.
\end{itemize}

%Все материалы размещены в открытом репозитории (лицензия позволяет свободное использование в образовательных целях). 
Комплекс готов к внедрению в учебный процесс; воспроизводимость окружения через Docker исключает проблемы с настройкой инструментов.